% Explicit second-order MF model and rotation products in the mixed family.
\section{混合試験族の明示的な二次MF復元とrotation積}
\label{kj:sec:second-jet}

本節では第\ref{mt:sec:family}節の非可換中心対について、二次の可換MF模型を計算する。
$h^3=0$で周囲の代数とKontsevich星積の同型および指定中心元の像を求め、
定理\ref{cg:thm:mf-comparison}によって内部シフト付きMF同値を得る。
MF pairの次数一・二のrotation積も決定する。
第\ref{am:sec:comparison}節の相対windowにより、
滑らかな古典fiberを持つ線束生成射影残余pairにも適用できる。

二次星積にはDitoの明示式\cite[式~(3)、arXiv版の印刷頁5]{KJDito}を使い、
曲率の量子化にはShoikhetのwheel消滅\cite[Section~1.3]{KJShoikhet}を使う。
これらは異なるgraphの重みである。
混合族への適用と明示的積表の先行研究上の初出は未確認であり、
formalityやgraph積分自体を本稿の新しい結果とは扱わない。

\subsection{星積の規格化とWeyl対称化}

$S=\CC[x,y,z,u,v]$、$R=\CC[\lambda,h]/(h^3)$とする。
$F=F_\lambda$を式\eqref{mt:eq:classical}の三次式として
\begin{equation}\label{kj:eq:pi}
 \pi=F_x\partial_y\wedge\partial_z+
 F_y\partial_z\wedge\partial_x+F_z\partial_x\wedge\partial_y
 =\frac12\pi^{ij}\partial_i\wedge\partial_j
\end{equation}
と置く。添字は五変数を走り、$u,v$を含む$\pi^{ij}$は零である。
$\{x_i,x_j\}=\pi^{ij}$という規格化で、$F,u,v$はCasimirである。
角度型Kontsevich星積は
\begin{align}
 f*g={}&fg+\frac h2\pi^{ij}f_i g_j\notag\\
 &+h^2\left\{\frac18\pi^{ij}\pi^{kl}f_{ik}g_{jl}
 +\frac1{12}\pi^{ij}\partial_j\pi^{kl}(f_{ik}g_l-f_k g_{il})
 -\frac1{24}Q_{ik}f_i g_k\right\},\label{kj:eq:star}\\
 Q_{ik}:={}&\sum_{j,l}(\partial_l\pi^{ij})(\partial_j\pi^{kl}).
 \label{kj:eq:eye-tensor}
\end{align}
同じ添字を足し合わせる。
Ditoの一次項は$\hbar\pi_{\rm Dito}^{ij}f_i g_j$なので、
$\hbar=h$、$\pi_{\rm Dito}=\pi/2$とした。
その二次式の最後の係数$-1/6$は本稿では$-1/24$になる。
この最後の項をeye項と呼ぶ。二次作用素全体は$f,g$について対称である。

星積によるWeyl対称化を
\[
 T(x_{i_1}\cdots x_{i_d})
 =\frac1{d!}\sum_{\sigma\in\mathfrak S_d}
 x_{i_{\sigma(1)}}*\cdots*x_{i_{\sigma(d)}}
\]
で定める。
\begin{lemma}[対称化写像の二次項]\label{kj:lem:symmetrization}
任意の多項式$f$について
\begin{equation}\label{kj:eq:T}
 T(f)=f-\frac{h^2}{48}\sum_{i,k}Q_{ik}\partial_i\partial_kf
\end{equation}
が成り立つ。これはJacobianイデアルを法としない多項式の等式であり、
任意の変数の数と一般のPoisson tensorについて有効である。
\end{lemma}
\begin{proof}
定数係数の一次式$\ell$について$B_1(\ell^m,\ell)=0$である。
二次項では$\ell$の二階微分を含む項は消え、
残るtree項も$\partial_j\pi^{kl}\ell_k\ell_l=0$で消える。
従って
\[
 B_2(\ell^m,\ell)=-\frac m{24}Q_{ik}\ell_i\ell_k\ell^{m-1}.
\]
既に$h^2$を含む補正に一次星積を作用させても$h^3$以上にしか寄与しない。
帰納法により
\[
 T(\ell^m)=\ell^m-\frac{h^2m(m-1)}{48}
 Q_{ik}\ell_i\ell_k\ell^{m-2}.
\]
各斉次次数の多項式は一次式の冪で張られるから、偏極により主張を得る。
\end{proof}

active Hessianを
\begin{equation}\label{kj:eq:hessian}
 \mathsf H=(F_{ij})_{i,j\in\{x,y,z\}},\qquad
 \mathcal H=\det\mathsf H=8xyz-2\lambda^2u^2z
\end{equation}
とする。
\begin{lemma}[Jacobian tensorとHessian]\label{kj:lem:hessian}
$Q$のactive部分は$-2\operatorname{adj}(\mathsf H)$、他の成分は零であり、
\begin{equation}\label{kj:eq:T-F}
 T(F)=F+\frac{h^2}{8}\mathcal H
\end{equation}
となる。tensor恒等式は、passive変数を係数とする任意の三変数多項式にも成り立つ。
\end{lemma}
\begin{proof}
交代tensorを$\epsilon_{ijk}$とすると$\pi^{ij}=\epsilon_{ijr}F_r$である。
従って余因子の表示から
\[
 Q_{ik}=\epsilon_{ijr}\epsilon_{kls}\mathsf H_{rl}\mathsf H_{sj}
       =-2\operatorname{adj}(\mathsf H)_{ik}.
\]
可逆性を仮定せず成立する恒等式である。$F_{ik}$と縮約すると
$\sum Q_{ik}F_{ik}=-6\det\mathsf H$となり、式\eqref{kj:eq:T}を使えばよい。
$u,v$を含む成分は$\pi$の脚がactive方向にしかないため零である。
\end{proof}

\subsection{指定された非可換中心元の像}

\begin{proposition}[周囲の代数と中心元の二次同定]
\label{kj:prop:algebra-and-center}
$A_R=A_{h,\lambda}/(h^3)$を式\eqref{mt:eq:relations}の代数とする。
生成元を同名の一次式へ送る写像は次数付き$R$-代数同型
\begin{equation}\label{kj:eq:Phi}
 \Phi:A_R\xrightarrow{\sim}(S\otimes R,*)
\end{equation}
を定め、式\eqref{mt:eq:center}の中心元について
\begin{equation}\label{kj:eq:effective}
 \begin{split}
 \Phi(\Omega_{h,\lambda})&=\widehat F_\lambda:=F_\lambda+h^2G_2,\\
 G_2&=\frac{\mathcal H}{24}
     =\frac{xyz}{3}-\frac{\lambda^2u^2z}{12}
 \end{split}
\end{equation}
となる。この命題は$\lambda^3\ne-3$を仮定しない。
\end{proposition}
\begin{proof}
二次作用素の対称性から$[x_i,x_j]_* =h\pi^{ij}$である。
関係式の右辺の二次式を星積で評価した際の補正には$h^2$が掛かり、
関係式の係数$h$と合わせると消える。例えば
\[
 [y,z]_* =h(x^2+\lambda uy)=h(x*x+\lambda u*y).
\]
他の関係も同様で、$u,v$との星積は通常の積である。
従って$\Phi$はwell-definedである。
補題\ref{mt:lem:artin-pbw}により各次数の定義域と値域は同階数の自由$R$-加群である。
$h=0$で恒等なので、冪零イデアル$(h)$について逆行列を有限展開でき、$\Phi$は同型となる。

通常の多項式から$A_R$へのWeyl対称化を$\sym_A$と書く。
中心元の閉公式は
\begin{equation}\label{kj:eq:omega-weyl}
 \Omega_{h,\lambda}=\sym_A(F_\lambda-h^2\mathcal H/12)
\end{equation}
と書ける。$\sym_A(xyz)$は六語の和の$1/6$倍で、
$\sym_A(\lambda uxy)=\lambda u(xy+yx)/2$である。
$\Phi\circ\sym_A=T$なので
\[
 \Phi(\Omega)=T(F)-\frac{h^2}{12}T(\mathcal H)
 =F+h^2\left(\frac18-\frac1{12}\right)\mathcal H
 =F+\frac{h^2}{24}\mathcal H .
\]
\end{proof}

中心元のWeyl記号$F-h^2\mathcal H/12$と、
Kontsevich側の曲率$F+h^2\mathcal H/24$は異なる。
差は対称化の$+h^2\mathcal H/8$であり、中心元の語を単に可換に置き換えても得られない。

\subsection{曲率のwheelと有限階数MF比較}

Kontsevichのaffine formality morphismを$\mathcal U$とし、
一次成分を通常の反対称化HKRに規格化する\cite[Sections~4--6]{PKontsevich}。
関数$f$に対するtangent mapは
\begin{equation}\label{kj:eq:tangent-map}
 \mathcal U^{h\pi}(f)=f+h\mathcal U_2(\pi,f)
 +\frac{h^2}{2}\mathcal U_3(\pi,\pi,f).
\end{equation}
\begin{lemma}[曲率側の二輪は消える]\label{kj:lem:curvature-wheel}
角度型formalityについて、任意の関数$f$に対し
$\mathcal U^{h\pi}(f)=f$が$h^3=0$で成り立つ。
\end{lemma}
\begin{proof}
出力は零項演算でありboundary vertexはない。
$\mathcal U_2(\pi,f)$では$\pi$からの二本の辺の終点は$f$しかなく、
自己辺と重複辺が許されないためgraphが存在しない。
$\mathcal U_3(\pi,\pi,f)$では、各$\pi$頂点から他方の$\pi$頂点と$f$へ
一本ずつ辺を出す二輪だけが可能である。
これは\cite[Section~8.3.3.1]{PKontsevich}のtangent-wheelで、
重み$W_2$は\cite[Section~1.3]{KJShoikhet}により零である。
重みはgraphだけで決まる定数なので$\pi$が線形である必要はない。
\end{proof}
星積のeye項には二つのboundary vertexがあり、その重みは零ではない。
曲率のgraphでは$f$がinternal vertexに置かれ、boundary vertexは零個である。
同じ$\partial\pi\,\partial\pi$というtensorが現れることだけから重みを同一視してはいけない。

\begin{theorem}[混合試験族の二次MF復元]
\label{kj:thm:curved-MC}\label{kj:thm:mf-reconstruction}
$\lambda$を任意の複素数に固定し、$R_\lambda=\CC[h]/(h^3)$とする。
第\ref{cg:sec:comparison}節の有限階数・冪等完備なMF圏について、
$h=0$で恒等となる$R_\lambda$-linearなdg Morita同値
\begin{equation}\label{kj:eq:mf-reconstruction}
 \bigl(\MF^{\ZZ}(A_R,\Omega_{h,\lambda}),(1)\bigr)
 \simeq
 \bigl(\MF^{\ZZ}(S\otimes R_\lambda,\widehat F_\lambda),(1)\bigr)
\end{equation}
が存在し、指定された三次周期性$(3)\simeq[2]$も保つ。
\end{theorem}
\begin{proof}
命題\ref{kj:prop:algebra-and-center}が生成元を固定する$\Phi$と
$\Phi(\Omega)=\widehat F_\lambda$を与える。
曲率wheelの消滅により$C_h(F)=F$、
$C_h(F+h^2G_2)=F+h^2G_2$なので、系\ref{cg:cor:matching}の条件を満たす。
同系の有限curved gaugeとtwisted complexへの延長を適用すればよい。
bar同型とその逆のTaylor成分はarity一・三・五に限られ、有限rankと内部shiftを保存する。
従って中心性やHochschild類の一致だけからの推論ではない。
\end{proof}

\subsection{低次数rotation積の完全な二次表}

$\mathcal M_R=\MF^{\ZZ}(A_R,\Omega_{h,\lambda})$、$\rho_R=(1)$と置く。
\begin{corollary}[MF pairの低次数積]\label{kj:cor:conditional-Nat}
$\lambda^3\ne-3$とする。
$R_\lambda$-linearな次数一・二の導来自然変換と積は
\[
 V=R_\lambda\langle x,y,z,u,v\rangle,\qquad
 W=\bigl(R_\lambda[x,y,z,u,v]/(\partial\widehat F_\lambda)\bigr)_2
\]
の積で与えられる。核の基底は
\begin{align}
 q_x&=x^2+\lambda uy+\frac{h^2}{3}yz,&
 q_y&=y^2+\lambda ux+\frac{h^2}{3}xz,\notag\\
 q_z&=z^2+h^2\left(\frac{xy}{3}-\frac{\lambda^2u^2}{12}\right),&
 q_u&=3u^2+\lambda xy-\frac{h^2\lambda^2}{6}uz,\notag\\
 q_v&=3v^2.&&\label{kj:eq:gradient-kernel}
\end{align}
$W$の順序付き基底を
\begin{equation}\label{kj:eq:W-basis}
 (xy,xz,xu,xv,yz,yu,yv,zu,zv,uv)
\end{equation}
とすれば、異なる二文字の積は対応する基底vectorで、平方の積は
\begin{align}
 x\cdot x&=-\lambda uy-\frac{h^2}{3}yz,&
 y\cdot y&=-\lambda ux-\frac{h^2}{3}xz,\notag\\
 z\cdot z&=-h^2\left(\frac13+\frac{\lambda^3}{36}\right)xy,&
 u\cdot u&=-\frac{\lambda}{3}xy+\frac{h^2\lambda^2}{18}uz,\notag\\
 v\cdot v&=0.&&\label{kj:eq:multiplication-table}
\end{align}
これは$10\times15$の対称積行列を完全に指定する。
\end{corollary}
\begin{proof}
定理\ref{kj:thm:mf-reconstruction}により可換側で計算できる。
可換超曲面は滑らかな特殊fiberを持つ冪零変形で、相対的に滑らかである。
命題\ref{prop:ej-geometric-low-degree}の対角Koszul計算を相対Artin基底で行う。
偏微分列は特殊fiberで正則列であり、有限な$h$進filtrationを用いると
正のKoszul次数の完全性が持ち上がる。
五変数では非恒等sectorはrotation次数一・二のcohomological degree零には寄与しない。
従って$V,W$と積は上の相対Jacobianデータである。

核は$\widehat F$の直接微分である。
平方$\{x^2,y^2,z^2,u^2,v^2\}$の係数小行列の行列式は$9$なので、
これら五関係は平方を一意に消去し、残る十単項式が$W$の自由基底となる。
$z^2$の式に$u^2=-\lambda xy/3\pmod{h^2}$を代入すれば表示した$z\cdot z$を得る。
他の積も同様である。
\end{proof}

\begin{corollary}[線束生成射影残余pairの二次復元]
\label{kj:cor:conditional-Ku}
$\lambda^3\ne-3$とし、非可換側の線束生成残余pairを
第\ref{am:sec:comparison}節の$\mathcal K_{B_R}^{\rm lin}$で定める。
このとき
\[
 (\mathcal K_{B_R}^{\rm lin},\rho_R)\simeq
 (\Ku_{\rm perf}(V(\widehat F_\lambda)/R_\lambda),\rho_{\widehat F_\lambda})
\]
であり、上の低次数積表はこのpairの積でもある。
可換側は$\operatorname{Perf}(V(\widehat F_\lambda))$内の残余成分を意味する。
\end{corollary}
\begin{proof}
定理\ref{am:thm:comparison}を非可換側と可換側に適用し、
定理\ref{kj:thm:mf-reconstruction}と合成する。
可換側の線束生成圏とperfect圏の一致は
系\ref{am:cor:formal-residual}の証明で確認する。
\end{proof}
Artin底の全有界coherent圏を、perfect圏と同一視していない。

古典Jacobian環での二次項は
\begin{equation}\label{kj:eq:scalar}
 [G_2]=c(\lambda)[xyz],\qquad
 c(\lambda)=\frac13+\frac{\lambda^3}{36}.
\end{equation}
実際、多項式恒等式
\begin{equation}\label{kj:eq:scalar-certificate}
 G_2-\left(\frac13+\frac{\lambda^3}{36}\right)xyz
 =-\frac{\lambda^2z}{36}\partial_uF_\lambda
\end{equation}
が証明書である。
さらに$u'=u-h^2\lambda^2z/36$という線形座標変換により
\begin{equation}\label{kj:eq:marked-normal-form}
 \widehat F_\lambda(x,y,z,u,v)
 =F_\lambda(x,y,z,u',v)+h^2\left(\frac13+\frac{\lambda^3}{36}\right)xyz
 \pmod{h^3}.
\end{equation}
従って$\lambda^3=-12$では、このMF pairは二次まで元の可換cubicのpairと同値になる。
この点でも古典cubicは滑らかであるが、三次以上の変形の消滅は従わない。
$\lambda$自体は切断していない。
$\lambda=0$の$c(0)=1/3$との一致は独立な整合確認であり、
係数$1/36$はその特殊化から推定したものではなく二次graphの係数から求めた値である。

\subsection{厳密計算の再現範囲}

\texttt{computations/verify\_kontsevich\_second\_jet.py}は有理数上で、
一般対称$3\times3$行列の$Q=-2\operatorname{adj}(\mathsf H)$、
混合族の三関係、三次の全35単項式の星積対称化、
生成元の125通りの結合性、中心元の像、
$5\times15$勾配行列と$10\times15$積行列の核・像を検算する。
平方小行列の行列式$9$、Euler積分、式\eqref{kj:eq:scalar-certificate}も確認し、
指定された出力先へJSON証明書を再生成する。

これらは多項式恒等式の検算であり、formality、wheel積分、
有限curved gaugeからのMF同値は計算機検証していない。
積行列のMF解釈は本文の比較定理と相対sector計算に依存する。
有限計算の再現範囲と、組版および他の証明の検証範囲は
区別して記録する。
